%% resume_template.tex -- starting point for every tailored resume.
%% Copy, rename to  Lastname_Company_RoleShort_ReqID.tex , fill, compile with
%% pdfLaTeX. Needs resume.cls beside it (or on TEXINPUTS).
%%
%% Angle-bracket text is placeholder. Nothing in angle brackets may survive
%% into a delivered file. Every claim must trace to the career profile.
\documentclass{resume}

% ---------- page-fit knobs (defaults shown; tighten in this order) ----------
% \setlength{\resitemsep}{2pt}     % 2pt -> 1pt -> 0pt
% \setlength{\resrolegap}{7pt}     % 7pt -> 5pt
% \setlength{\ressecbefore}{10pt}  % 10pt -> 8pt
% Never touch font size or margins. After the knobs, cut the weakest bullet.

\name{<First Last>}
% Separate items with \sep. City is optional. Anything else (citizenship, an
% active clearance, a portfolio link) is added as one more \sep item, and only
% when the user's profile says to show it.
\contact{<City, ST> \sep \href{mailto:<email>}{<email>} \sep <(000) 000-0000>}

\begin{document}
\makeheader

% Three sentences, no first person, no years-of-experience count.
% 1) identity noun phrase that mirrors the posting's discipline
% 2) what the candidate delivers, in the posting's vocabulary
% 3) the differentiating background
\section*{Professional Summary}
<Summary.>

% 5 to 7 lines. Rename and reorder the categories for each posting so line 1
% is the posting's center of gravity. Every item needs evidence in a bullet
% below or a confirmed line in the career profile.
\section*{Technical Skills}
\skill{<Category>}{<Item, Item, Item>}
\skill{<Category>}{<Item, Item, Item>}

% Reverse chronological, always. Shift emphasis with bullet count and bullet
% order, never by reordering jobs. Titles are verbatim from the profile.
\section*{Professional Experience}
\role{<Official Title>}{<Month YYYY> -- Present}
     {<Employer (context if needed)> -- <City, ST>}
\begin{itemize}
  \item <Action verb + what + how (tools, standards) + result. Present tense for ongoing duties, past tense for finished projects.>
  \item <Most posting-relevant bullets first.>
\end{itemize}

% Put \pagebreak between roles when a role would otherwise split across pages.
% \pagebreak

\role{<Official Title>}{<Month YYYY> -- <Month YYYY>}
     {<Employer> -- <City, ST>}
\begin{itemize}
  \item <Past tense.>
\end{itemize}

\section*{Education}
\degree{<M.S. Field, GPA: 0.00>}{<Month YYYY>}
       {<University>}
       {<Thesis: Title>}
\degree{<B.S. Field, GPA: 0.00 (Honors)>}{<Month YYYY>}
       {<University>}
       {}

% Use "Publications" only if every entry is a paper. If any entry is a
% poster, talk, or report, title the section "Publications \& Presentations"
% and tag that entry, e.g. (Poster).
\section*{Publications \& Presentations}
\publication{<Title>}
            {<Venue, Year>}

% One form for the whole section: every project followed by a plain line, or
% every project followed by an itemize. Never mix.
\section*{Projects}
\project{<Project Title>}{<Context or award>}
<One-line description.>

\end{document}
